% --- Archon physics formalization source begin ---
% archon:physics
% archon:covers PhyXMiniProblems/problem_phyx_mini_0550.lean
% archon:source-report reports/phyx_mini/problem_phyx_mini_0550.source.json

\chapter{Physics problem phyx\_mini\_0550}
\label{ch:PhyXMiniProblems_problem_phyx_mini_0550}

\paragraph{Problem source.}
\#\# Physical scenario

A particle is confined between two impenetrable walls separated by a distance \$l\$, as shown in the figure. The wave function describing the particle's state is \$\textbackslash{}psi = cx(l-x)\$, where \$c\$ is an undetermined constant.

\#\# Question

Calculate the probability of finding the particle in the interval from \$0\$ to \$\textbackslash{}frac\{1\}\{3\}l\$.

\#\# Answer choices

- A: A: \textbackslash{}[ \textbackslash{}frac\{11\}\{72\} \textbackslash{}]

- B: B: \textbackslash{}[ \textbackslash{}frac\{35\}\{72\} \textbackslash{}]

- C: C: \textbackslash{}[ \textbackslash{}frac\{37\}\{81\} \textbackslash{}]

- D: D: \textbackslash{}[ \textbackslash{}frac\{17\}\{81\} \textbackslash{}]

\#\# Figure caption (auxiliary; use the image as primary evidence)

The image depicts a diagram featuring a horizontal line with an arrow pointing to the right, labeled with an "x" at the end, indicating the positive x-direction. Along this horizontal line, there are two vertical, hatched lines.

- The left vertical line is at the origin, labeled "O."
- The right vertical line is at a distance "l" from the origin along the horizontal line.
- Approximately one-third of the way from the origin "O" to the point labeled "l," there is a mark with the label "1/3 l."

The diagram appears to be representing a spatial relationship between two points on a line, possibly for geometric or physical analysis.

\#\# Dataset metadata

Quantum Phenomena; Numerical Reasoning, Physical Model Grounding Reasoning

\paragraph{Recorded answer/context.}
D: D: \textbackslash{}[ \textbackslash{}frac\{17\}\{81\} \textbackslash{}]

\paragraph{Figure/image path.}
/root/proposal\_for\_physic/science-mango/phyx\_mini\_run/phyx\_data/test\_image/550.png

\paragraph{Formalization target.}
create a compiling Lean file with sorry bodies at `PhyXMiniProblems/problem\_phyx\_mini\_0550.lean`. The Lean declarations must preserve the physical quantities, dimensions or dimensional roles, figure labels, governing-law hypotheses, and final relation expressed by this problem.
Use Mathlib/Physlib names found through LeanExplore where available. If a physics API is missing, introduce faithful local abstractions rather than scalar placeholder aliases.

\begin{theorem}[Physics formalization target]
\label{thm:physics:phyx_mini_0550:target}
\lean{PhyXMiniProblems.ProblemPhyXMini0550.probability_from_origin_to_one_third}
\uses{def:physics:phyx-mini-0550:phyxminiproblems-problemphyxmini0550-particleinboxsetup, def:physics:phyx-mini-0550:phyxminiproblems-problemphyxmini0550-matchessuppliedparticleinboxfigure, def:physics:phyx-mini-0550:phyxminiproblems-problemphyxmini0550-hasquadraticwaveprofile, def:physics:phyx-mini-0550:phyxminiproblems-problemphyxmini0550-satisfiesimpenetrablewallboundarycondition, def:physics:phyx-mini-0550:phyxminiproblems-problemphyxmini0550-issquareintegrablepositionstate, def:physics:phyx-mini-0550:phyxminiproblems-problemphyxmini0550-isnormalizedinsidebox, def:physics:phyx-mini-0550:phyxminiproblems-problemphyxmini0550-satisfiespositionbornrule}
The assigned autoformalize agent should translate this physics problem into Lean declarations in the covered file, with theorem and lemma proof bodies written as `by sorry`.
\end{theorem}
\begin{proof}
This is an autoformalization task, not a proof task. Produce statements that can later be proved by the physics prover without weakening the physical model.
\end{proof}

% --- Archon declaration topology begin ---
\section{Declaration topology}
\begin{definition}[Position Quantity]
  \label{def:physics:phyx-mini-0550:phyxminiproblems-problemphyxmini0550-positionquantity}
  \lean{PhyXMiniProblems.ProblemPhyXMini0550.PositionQuantity}
  A signed physical position on the horizontal axis
\end{definition}
\begin{proof}
  This is the declared model definition of Position Quantity.
\end{proof}
\begin{definition}[Length Quantity]
  \label{def:physics:phyx-mini-0550:phyxminiproblems-problemphyxmini0550-lengthquantity}
  \lean{PhyXMiniProblems.ProblemPhyXMini0550.LengthQuantity}
  A nonnegative physical length, used for the wall separation l
\end{definition}
\begin{proof}
  This is the declared model definition of Length Quantity.
\end{proof}
\begin{definition}[Wave Amplitude Quantity]
  \label{def:physics:phyx-mini-0550:phyxminiproblems-problemphyxmini0550-waveamplitudequantity}
  \lean{PhyXMiniProblems.ProblemPhyXMini0550.WaveAmplitudeQuantity}
  A one-dimensional position-wavefunction amplitude, of dimension L⁻¹ᐟ²
\end{definition}
\begin{proof}
  This is the declared model definition of Wave Amplitude Quantity.
\end{proof}
\begin{definition}[Quadratic Coefficient Quantity]
  \label{def:physics:phyx-mini-0550:phyxminiproblems-problemphyxmini0550-quadraticcoefficientquantity}
  \lean{PhyXMiniProblems.ProblemPhyXMini0550.QuadraticCoefficientQuantity}
  The coefficient in ψ(x) = c x (l-x), of dimension L⁻⁵ᐟ²
\end{definition}
\begin{proof}
  This is the declared model definition of Quadratic Coefficient Quantity.
\end{proof}
\begin{definition}[position Readout]
  \label{def:physics:phyx-mini-0550:phyxminiproblems-problemphyxmini0550-positionreadout}
  \lean{PhyXMiniProblems.ProblemPhyXMini0550.positionReadout}
  \uses{def:physics:phyx-mini-0550:phyxminiproblems-problemphyxmini0550-positionquantity}
  Numerical position coordinate in a selected physical length unit
\end{definition}
\begin{proof}
  This is the declared model definition of position Readout.
\end{proof}
\begin{definition}[length Readout]
  \label{def:physics:phyx-mini-0550:phyxminiproblems-problemphyxmini0550-lengthreadout}
  \lean{PhyXMiniProblems.ProblemPhyXMini0550.lengthReadout}
  \uses{def:physics:phyx-mini-0550:phyxminiproblems-problemphyxmini0550-lengthquantity}
  Numerical nonnegative length in a selected physical length unit
\end{definition}
\begin{proof}
  This is the declared model definition of length Readout.
\end{proof}
\begin{definition}[position From Readout]
  \label{def:physics:phyx-mini-0550:phyxminiproblems-problemphyxmini0550-positionfromreadout}
  \lean{PhyXMiniProblems.ProblemPhyXMini0550.positionFromReadout}
  \uses{def:physics:phyx-mini-0550:phyxminiproblems-problemphyxmini0550-positionquantity}
  Construct a physical position from its coordinate in a selected unit
\end{definition}
\begin{proof}
  This is the declared model definition of position From Readout.
\end{proof}
\begin{definition}[wave Amplitude Readout]
  \label{def:physics:phyx-mini-0550:phyxminiproblems-problemphyxmini0550-waveamplitudereadout}
  \lean{PhyXMiniProblems.ProblemPhyXMini0550.waveAmplitudeReadout}
  \uses{def:physics:phyx-mini-0550:phyxminiproblems-problemphyxmini0550-waveamplitudequantity}
  Numerical wave-amplitude readout in inverse square-root selected units
\end{definition}
\begin{proof}
  This is the declared model definition of wave Amplitude Readout.
\end{proof}
\begin{definition}[coefficient Readout]
  \label{def:physics:phyx-mini-0550:phyxminiproblems-problemphyxmini0550-coefficientreadout}
  \lean{PhyXMiniProblems.ProblemPhyXMini0550.coefficientReadout}
  \uses{def:physics:phyx-mini-0550:phyxminiproblems-problemphyxmini0550-quadraticcoefficientquantity}
  Numerical readout of c in inverse five-halves powers of the selected unit
\end{definition}
\begin{proof}
  This is the declared model definition of coefficient Readout.
\end{proof}
\begin{definition}[Wall Rendering]
  \label{def:physics:phyx-mini-0550:phyxminiproblems-problemphyxmini0550-wallrendering}
  \lean{PhyXMiniProblems.ProblemPhyXMini0550.WallRendering}
  The visual style used for each vertical wall in the supplied diagram
\end{definition}
\begin{proof}
  This is the declared model definition of Wall Rendering.
\end{proof}
\begin{definition}[One Dimensional Position State]
  \label{def:physics:phyx-mini-0550:phyxminiproblems-problemphyxmini0550-onedimensionalpositionstate}
  \lean{PhyXMiniProblems.ProblemPhyXMini0550.OneDimensionalPositionState}
  \uses{def:physics:phyx-mini-0550:phyxminiproblems-problemphyxmini0550-positionquantity, def:physics:phyx-mini-0550:phyxminiproblems-problemphyxmini0550-waveamplitudequantity}
  A pointwise one-dimensional quantum state together with its dimensionless position-probability observable. Born's rule is imposed separately, so the observable is not defined using the requested numerical answer
\end{definition}
\begin{proof}
  This is the declared model definition of One Dimensional Position State.
\end{proof}
\begin{definition}[Particle In Box Figure]
  \label{def:physics:phyx-mini-0550:phyxminiproblems-problemphyxmini0550-particleinboxfigure}
  \lean{PhyXMiniProblems.ProblemPhyXMini0550.ParticleInBoxFigure}
  \uses{def:physics:phyx-mini-0550:phyxminiproblems-problemphyxmini0550-positionquantity, def:physics:phyx-mini-0550:phyxminiproblems-problemphyxmini0550-lengthquantity, def:physics:phyx-mini-0550:phyxminiproblems-problemphyxmini0550-wallrendering}
  Geometry, labels, and landmarks shown in the supplied primary figure
\end{definition}
\begin{proof}
  This is the declared model definition of Particle In Box Figure.
\end{proof}
\begin{definition}[Particle In Box Setup]
  \label{def:physics:phyx-mini-0550:phyxminiproblems-problemphyxmini0550-particleinboxsetup}
  \lean{PhyXMiniProblems.ProblemPhyXMini0550.ParticleInBoxSetup}
  \uses{def:physics:phyx-mini-0550:phyxminiproblems-problemphyxmini0550-quadraticcoefficientquantity, def:physics:phyx-mini-0550:phyxminiproblems-problemphyxmini0550-onedimensionalpositionstate, def:physics:phyx-mini-0550:phyxminiproblems-problemphyxmini0550-particleinboxfigure}
  The particle state and the undetermined coefficient appearing in its profile
\end{definition}
\begin{proof}
  This is the declared model definition of Particle In Box Setup.
\end{proof}
\begin{definition}[Matches Supplied Particle In Box Figure]
  \label{def:physics:phyx-mini-0550:phyxminiproblems-problemphyxmini0550-matchessuppliedparticleinboxfigure}
  \lean{PhyXMiniProblems.ProblemPhyXMini0550.MatchesSuppliedParticleInBoxFigure}
  \uses{def:physics:phyx-mini-0550:phyxminiproblems-problemphyxmini0550-positionreadout, def:physics:phyx-mini-0550:phyxminiproblems-problemphyxmini0550-lengthreadout, def:physics:phyx-mini-0550:phyxminiproblems-problemphyxmini0550-particleinboxfigure}
  Exact labels, wall styles, orientation, and relative positions read from the primary image. These are figure facts; no probability value occurs here
\end{definition}
\begin{proof}
  This is the declared model definition of Matches Supplied Particle In Box Figure.
\end{proof}
\begin{definition}[Has Quadratic Wave Profile]
  \label{def:physics:phyx-mini-0550:phyxminiproblems-problemphyxmini0550-hasquadraticwaveprofile}
  \lean{PhyXMiniProblems.ProblemPhyXMini0550.HasQuadraticWaveProfile}
  \uses{def:physics:phyx-mini-0550:phyxminiproblems-problemphyxmini0550-positionreadout, def:physics:phyx-mini-0550:phyxminiproblems-problemphyxmini0550-positionfromreadout, def:physics:phyx-mini-0550:phyxminiproblems-problemphyxmini0550-waveamplitudereadout, def:physics:phyx-mini-0550:phyxminiproblems-problemphyxmini0550-coefficientreadout, def:physics:phyx-mini-0550:phyxminiproblems-problemphyxmini0550-particleinboxsetup}
  Inside the walls, the amplitude has the stated profile ψ(x) = c (x-O) (l-(x-O)), equivalently c (x-O) (x\_right-x)
\end{definition}
\begin{proof}
  This is the declared model definition of Has Quadratic Wave Profile.
\end{proof}
\begin{definition}[Satisfies Impenetrable Wall Boundary Condition]
  \label{def:physics:phyx-mini-0550:phyxminiproblems-problemphyxmini0550-satisfiesimpenetrablewallboundarycondition}
  \lean{PhyXMiniProblems.ProblemPhyXMini0550.SatisfiesImpenetrableWallBoundaryCondition}
  \uses{def:physics:phyx-mini-0550:phyxminiproblems-problemphyxmini0550-positionreadout, def:physics:phyx-mini-0550:phyxminiproblems-problemphyxmini0550-positionfromreadout, def:physics:phyx-mini-0550:phyxminiproblems-problemphyxmini0550-waveamplitudereadout, def:physics:phyx-mini-0550:phyxminiproblems-problemphyxmini0550-particleinboxsetup}
  Impenetrable-wall confinement: the position amplitude vanishes outside the box
\end{definition}
\begin{proof}
  This is the declared model definition of Satisfies Impenetrable Wall Boundary Condition.
\end{proof}
\begin{definition}[Is Square Integrable Position State]
  \label{def:physics:phyx-mini-0550:phyxminiproblems-problemphyxmini0550-issquareintegrablepositionstate}
  \lean{PhyXMiniProblems.ProblemPhyXMini0550.IsSquareIntegrablePositionState}
  \uses{def:physics:phyx-mini-0550:phyxminiproblems-problemphyxmini0550-positionfromreadout, def:physics:phyx-mini-0550:phyxminiproblems-problemphyxmini0550-waveamplitudereadout, def:physics:phyx-mini-0550:phyxminiproblems-problemphyxmini0550-particleinboxsetup}
  The pointwise representative of the state is square-integrable in every unit
\end{definition}
\begin{proof}
  This is the declared model definition of Is Square Integrable Position State.
\end{proof}
\begin{definition}[Is Normalized Inside Box]
  \label{def:physics:phyx-mini-0550:phyxminiproblems-problemphyxmini0550-isnormalizedinsidebox}
  \lean{PhyXMiniProblems.ProblemPhyXMini0550.IsNormalizedInsideBox}
  \uses{def:physics:phyx-mini-0550:phyxminiproblems-problemphyxmini0550-positionreadout, def:physics:phyx-mini-0550:phyxminiproblems-problemphyxmini0550-positionfromreadout, def:physics:phyx-mini-0550:phyxminiproblems-problemphyxmini0550-waveamplitudereadout, def:physics:phyx-mini-0550:phyxminiproblems-problemphyxmini0550-particleinboxsetup}
  The confined wavefunction has total probability one inside the two walls
\end{definition}
\begin{proof}
  This is the declared model definition of Is Normalized Inside Box.
\end{proof}
\begin{definition}[Satisfies Position Born Rule]
  \label{def:physics:phyx-mini-0550:phyxminiproblems-problemphyxmini0550-satisfiespositionbornrule}
  \lean{PhyXMiniProblems.ProblemPhyXMini0550.SatisfiesPositionBornRule}
  \uses{def:physics:phyx-mini-0550:phyxminiproblems-problemphyxmini0550-positionquantity, def:physics:phyx-mini-0550:phyxminiproblems-problemphyxmini0550-positionreadout, def:physics:phyx-mini-0550:phyxminiproblems-problemphyxmini0550-positionfromreadout, def:physics:phyx-mini-0550:phyxminiproblems-problemphyxmini0550-waveamplitudereadout, def:physics:phyx-mini-0550:phyxminiproblems-problemphyxmini0550-particleinboxsetup}
  The position-space Born rule: probability between any ordered pair of physical positions is the integral of the squared complex amplitude over that interval
\end{definition}
\begin{proof}
  This is the declared model definition of Satisfies Position Born Rule.
\end{proof}
% --- Archon declaration topology end ---
% --- Archon physics formalization source end ---
